\documentclass[12pt]{article}
% \usepackage[top=1in,left=1in, right = 1in, footskip=1in]{geometry}
\usepackage[top=1in,footskip=1in]{geometry}

\usepackage{graphicx}
\usepackage{xspace}
%\usepackage{adjustbox}

\usepackage{multirow}
\usepackage{booktabs}

\usepackage{pdflscape}

\usepackage{grffile}

\newcommand{\comment}{\showcomment}
%% \newcommand{\comment}{\nocomment}

\newcommand{\showcomment}[3]{\textcolor{#1}{\textbf{[#2: }\textsl{#3}\textbf{]}}}
\newcommand{\nocomment}[3]{}

\newcommand{\jd}[1]{\comment{cyan}{JD}{#1}}
\newcommand{\swp}[1]{\comment{magenta}{SWP}{#1}}
\newcommand{\bmb}[1]{\comment{blue}{BMB}{#1}}
\newcommand{\djde}[1]{\comment{red}{DJDE}{#1}}

\newcommand{\eref}[1]{Eq.~(\ref{eq:#1})}
\newcommand{\fref}[1]{Fig.~\ref{fig:#1}}
\newcommand{\Fref}[1]{Fig.~\ref{fig:#1}}
\newcommand{\sref}[1]{Sec.~\ref{#1}}
\newcommand{\frange}[2]{Fig.~\ref{fig:#1}--\ref{fig:#2}}
\newcommand{\tref}[1]{Table~\ref{tab:#1}}
\newcommand{\tlab}[1]{\label{tab:#1}}
\newcommand{\seminar}{SE\mbox{$^m$}I\mbox{$^n$}R}

\usepackage{amsthm}
\usepackage{amsmath}
\usepackage{amssymb}
\usepackage{amsfonts}
\usepackage[utf8]{inputenc} % make sure fancy dashes etc. don't get dropped

\usepackage{booktabs}
\usepackage{tabularx}
\usepackage{array}

\newcolumntype{L}[1]{>{\raggedright\arraybackslash}p{#1}}
\newcolumntype{C}[1]{>{\centering\arraybackslash}p{#1}}
\newcolumntype{Y}{>{\raggedright\arraybackslash}X}

\usepackage{lineno}
\linenumbers

\usepackage[pdfencoding=auto, psdextra]{hyperref}

\usepackage{natbib}
\bibliographystyle{unsrt}
\date{\today}

\usepackage{xspace}
\newcommand*{\ie}{i.e.\@\xspace}

\usepackage{color}

\newcommand{\Rx}[1]{\ensuremath{{\mathcal R}_{#1}}\xspace} 
\newcommand{\RR}{\ensuremath{{\mathcal R}}\xspace}
\newcommand{\Rres}{\Rx{\mathrm{res}}}
\newcommand{\Rinv}{\Rx{\mathrm{inv}}}
\newcommand{\Rhat}{\ensuremath{{\hat\RR}}}
\newcommand{\Rt}{\ensuremath{{\mathcal R}(t)}\xspace}
\newcommand{\tsub}[2]{#1_{{\textrm{\tiny #2}}}}
\newcommand{\dd}[1]{\ensuremath{\, \mathrm{d}#1}}
\newcommand{\dtau}{\dd{\tau}}
\newcommand{\dx}{\dd{x}}
\newcommand{\dsigma}{\dd{\sigma}}

\newcommand{\rx}[1]{\ensuremath{{r}_{#1}}\xspace} 
\newcommand{\rres}{\rx{\mathrm{res}}}
\newcommand{\rinv}{\rx{\mathrm{inv}}}

\newcommand{\psymp}{\ensuremath{p}} %% primary symptom time
\newcommand{\ssymp}{\ensuremath{s}} %% secondary symptom time
\newcommand{\pinf}{\ensuremath{\alpha_1}} %% primary infection time
\newcommand{\sinf}{\ensuremath{\alpha_2}} %% secondary infection time

\newcommand{\psize}{{\mathcal P}} %% primary cohort size
\newcommand{\ssize}{{\mathcal S}} %% secondary cohort size

\newcommand{\gtime}{\tau_{\rm g}} %% generation interval
\newcommand{\gdist}{g} %% generation-interval distribution
\newcommand{\idist}{\ell} %% incubation-period distribution

\newcommand{\total}{{\mathcal T}} %% total number of serial intervals

\usepackage{lettrine}

\newcommand{\dropcapfont}{\fontfamily{lmss}\bfseries\fontsize{26pt}{28pt}\selectfont}
\newcommand{\dropcap}[1]{\lettrine[lines=2,lraise=0.05,findent=0.1em, nindent=0em]{{\dropcapfont{#1}}}{}}

\begin{document}

\begin{flushleft}{
	\Large
	\textbf\newline{
	  Personalized AI assistants could either promote or hinder the emergence of collective intelligence
	}
}
\newline
\\
Jaeseung Kim\textsuperscript{1} Sang Woo Park\textsuperscript{1,2,3,*}
\\
\bigskip
\textbf{1} School of Biological Sciences, Seoul National University, Seoul, Korea
\bigskip

*Corresponding author: sangwoopark@snu.ac.kr
\end{flushleft}

\section*{Abstract}

Advances in artificial intelligence (AI) have enabled the wide use of personalized assistance for cognitive tasks, which enhanced individual performance on various fields.
However, the impact of AI assistance on collective intelligence remains poorly understood. AI assistance may undermine the emergence of collective intelligence by reducing the individual diversity or introducing structured biases into a population. 
Here, we develop a human-AI collective prediction model mediated by social learning based on incentive structure by extending the existing works on collective intelligence.
Our model uncovers that emergence of collective intelligence is depend on how human and AI interact. 
When AI works as if it is omniscient and gives assistance regardless of users' interests, collective intelligence will be promoted.
However, when AI works as a chat-bot and gives assistance only based on users' interests, collective intelligence could be hindered due to diversity loss and overconfidence on AI.
We showed that, to promote collective intelligence, penalizing or incentivizing AI usage is not reliable. Instead, elaborating more diversity-aware incentives could be the key to promote and maintain the collective intelligence.
\pagebreak

\section*{Introduction}

Collective intelligence (CI) describes a phenomenon in which a group’s ability to solve problems, make decisions, or adapt to changing environments exceeds each member’s ability \cite{galton1907vox,woolley2010evidence,yaniv2007using,hommes2005coordination,lorge1958survey}.
This phenomenon, also referred to as wisdom of crowds, has thus served as a central mechanism underlying the success of diverse organizations, ranging from animal groups \cite{dreyer2025comparing,simons2004many,berdahl2013emergent} and human societies \cite{woolley2010evidence, riedl2021quantifying} to machine-learning systems \cite{breiman1996bagging,breiman2001random}.
Understanding the conditions that promote and  maintain CI therefore remains a fundamental question across social sciences \cite{mann2017optimal}, evolutionary biology \cite{wang2025individual}, and computer science \cite{werfel2026fluid}.
Addressing this question has become particularly timely as artificial intelligence (AI) has begun to reshape how individuals acquire information, make decisions, and contribute to collective outcomes \cite{cui2024ai}.

The recent emergence of generative AIs and large language models (LLMs) has changed the way individuals engage in various cognitive tasks, ranging from simple everyday tasks to more complex tasks, including scientific discovery \cite{lu2026towards} and creative writing \cite{doshi2024generative}.
In particular, while the use of personal AI models, such as ChatGPT, has been rapidly increasing \cite{sajadieh2026artificial}, the potential long-term impact of personal AI services on human society, especially the CI, remains uncertain.
Some argue that AI can promote CI by enhancing individual expertise \cite{vaccaro2024combinations}, and opinion aggregation \cite{burton2024large}.
Conversely, empirical evidence indicate that the use of AIs can limit the diversity of opinions \cite{doshi2024generative,meincke2025chatgpt, fugener2021will}, a key requirement for the emergence of CI \cite{lorenz2011social,aplin2014individual,hong2004groups}.

Our ability to predict the long-term impact of personal AI services on human CI is hampered by the apparent discrepancy in the incentives that shape AI uses.
On one hand, many industrial sectors are pushing towards increased AI usages, where the number of ``tokens'' used is often perceived as a proxy for productivity \cite{bousquette2026tokenmaxxing,keary2026cult}.
On the other hand, many academic institutions appear to discourage, and even penalize, AI use, reflecting concerns about plagiarism, misinformation, and limited independnce \cite{chawla2026researchers,thorp2023chatgpt}.
Such discrepancy raises a question about whether directly incentivizing or penalizing AI usage will benefit or harm human CI.
These observations further highlight the importance of indentifying a robust incentive structure that can promote CI in the presence of AI, which can retain the benefits of AIs while preventing potential risk such as preserving the diversity of individual opinions.

To address these questions, we extend previous approaches for modeling CI, where each individual makes independent predictions for the outcome from partial observations of underlying causal factors \cite{mann2017optimal,wang2025individual,hong2012incentives}.
We allow each individual to independently consult with personal AIs and revise their predictions accordingly.
Across different assumptions about the AI model's ability and underlying incentive structures, we investigate the robustness of human CIs.

\section*{Main}
\subsection*{Model}

\swp{unclear what is meant by decentralized manner.}
\swp{how about something like this: We begin by extending previously established models of CI in which a group of players participates in a prediction task from partial/independent observations of underlying causal factors and aggregates their predictions...}
We begin by extending previously established models of CI in which a group of players makes collective predictions by gathering information and aggregating individual predictions in a decentralized manner \cite{hong2012incentives, mann2017optimal,wang2025individual}.
\swp{I would first describe Y and then each individual's prediction. I think we can also leave out ``in the previous approach''.}
In the previous approach, among $M$ different factors $(x_i)$ that affect the outcome $(Y)$, an individual ($A$) makes a prediction based on their factor of interest $(x_{e_A})$ and their own belief $(c_A)$: $Y_{A} = c_A x_{e_A}$.
Here, we assume each individual can independently consult AI and revise their prediction based on their reliance on AI ($\beta_A$): $Y_A^{re} = \beta_A Y_A^{AI} + (1 - \beta_A) Y_A$.
\swp{We need to say that this revised prediction now contributes to the collective outcome.}

\swp{I would first describe what AI does biologically. Noise and bias are secondary information and can be added towards the end.
For example: To model AI prediction, we consider a scenario where AI has already been trained on this prediction task and therefore has information about the effect of each causal factor.
Every time a player consults with AI, it will make predictions based on the known information, which can include factor-specific bias and stochastic error.
}
To specify how individuals consult AI, we model the AI prediction $(Y_A^{AI})$ as a noisy and biased estimate of the target outcome.
Specifically, AI predictions include factor-specific bias $(b_i)$ and stochastic error $(\epsilon)$.
We then consider two extreme modes of AI response, which differ in the amount of information AI uses when consulting an individual user \swp{when making prediction?}.
The first is the ``AI knows all'' model, in which AI predicts the outcome based on all factors regardless of the user's interest:
$Y_{A}^{AI} = \sum_{i = 0}^{M}(\alpha_{i}x_{i} + b_{i}) + \epsilon$
, where $\alpha_i$ represents the true coefficient of the factor.
The second is the ``AI answers question'' model, in which AI predicts the outcome based only on the user's interest:
$Y_{A}^{AI} = \alpha_{e_A}x_{e_A} + b_{e_{A}} + \epsilon$.

\swp{After every round of prediction task?}
After the prediction, each individual receives payoff under the assumed incentive structure.
We considere two incentive structures, ``feedback'' and ``niche expert'' \cite{wang2025individual}, which have been shown to promote near-perfect collective accuracy in the absence of AI (except averaging aggregation with niche expert incentive).
The feedback structure gives higher payoff to individuals who can push the collective prediction toward the true outcome, regardless of their individual predictive accuracy.
In contrast, the niche expert structure gives higher payoff to individuals who have both better predictive accuracy and a minority interest.

\swp{Need to say something about each round vs each generation... I think one generation consists of several rounds but we're making some kind of equilibrium assumption?}
Through time, the group evolves by imitating the peers who got better payoff.
Specifically, for each generation, two individuals are randomly sampled, and one imitates all three strategies (factor of interest, belief, reliance on AI) of the other with a probability determined by the difference in expected payoff.
Individual predictions are aggregated into a collective prediction by simple averaging in the AI knows all model, or by clustering in the AI answers question model, depending on the mode of AI.

\swp{I think we need a brief paragraph about what our goals are to make the next sections flow better.
For example: this model now allows us to explore the impact of AI on }

\subsection*{AI can either promote or prevent the emergence of CI}
To understand how AI could affect the emergence of CI, we first compare collective dynamics with and without AI assistance.
We find that AI affects CI in qualitatively different ways depending on the mode of AI.

In the AI knows all model, AI promotes CI by increasing either the speed at which CI emerges or the resulting collective accuracy (Fig. 2A).
When the group can achieve near-perfect CI without AI (feedback incentive), AI accelerates its emergence by raising the initial collective accuracy to the level of AI accuracy.
By contrast, when the group cannot achieve CI without AI (niche expert incentive), AI enables moderately accurate CI that exceeds the accuracy achievable by either AI or the human group alone.
These trends are robust across different levels of AI accuracy and can be shown mathematically by assuming a sufficiently large group of players (see SI).


In contrast, In the AI answers question model, AI prevents the emergence and maintenance of CI (Fig. 2B).
Under the feedback incentive, the group initially achieves near-perfect CI. 
However, as interest diversity rapidly decreases, the group fails to capture variation in some factors, leading to increased collective variance and a breakdown of collective accuracy.
Under the niche expert incentive, collective accuracy initially increases but peaks before near-perfect CI is achieved, and then plateaus with a slight decline.
This trend is mediated by increased reliance on AI: individuals with higher reliance on AI obtain higher payoffs, causing the group to become increasingly reliant on AI and the collective prediction to inherit AI bias.
However, the niche expert incentive can also drive reliance on AI toward zero, depending on the relationship between the true coefficient and AI bias.
In this case, the group excludes AI assistance and builds CI on its own, which largely delays the emergence of CI (see SI).

\subsection*{Directly incentivizing or penalizing AI use can harm CI}
Given the strong influence of AI use on collective outcomes, a straightforward intervention is to directly incentivize or penalize AI use itself.
To investigate whether such direct incentives or penalties improve CI, we modify the two incentive structures by adding a payoff term proportional to individual reliance on AI:
\begin{equation}
\mathrm{Feedback /} \mathrm{Niche\ expert} + \lambda \beta_A
\end{equation}
where $\lambda$ represent the strength of incentive ($\lambda > 0$) or penalty ($\lambda < 0$).
Except in the AI knows all model under the feedback incentive, directly incentivizing or penalizing AI use fails to promote near-perfect CI.
Moreover, in most cases, directly incentivizing or penalizing leads to lower collective accuracy than the original incentive structures (Fig. 3).

When AI use is directly incentivized, reliance on AI evolves toward full reliance ($\beta_A = 1$), reducing collective accuracy as the collective prediction inherits AI bias. 
By contrast, when AI use is directly penalized, group evolves toward no reliance on AI ($\beta_A = 0$), where collective accuracy decreases with the strength of the penalty.
In some cases, when AI bias is relatively large and the penalty is weak, penalizing AI use can achieve higher collective accuracy than the original incentive structure.
However, in most cases, penalizing AI use imposes strong selection against AI reliance, pushing the group toward lower stationary collective accuracy.
Although the AI answers question model under the feedback incentive already shows stochastic breakdown of CI due to the loss of interest diversity, incentivizing or penalizing AI use further strengthens the breakdown.

% \subsection*{Balanced incentives mitigate AI-induced failures of CI}
% Focusing on the AI answers question model, in which both incentive structures fail to promote accurate CI, we investigate whether a better incentive structure can overcome the risks introduced by AI.
% Here, we propose a ``balanced'' incentive, defined as the simple average of the feedback and niche expert incentives.

% The core idea of the balanced incentive is to maintain both within- and between-interest diversity.
% The feedback incentive can strongly maintain within-interest diversity, defined as diversity in strategies (belief and reliance on AI), among individuals with the same factor of interest.
% However, between-interest diversity, which we denoted as ``interest diversity'' in the previous sections, can be hindered when AI assistance is present.
% This occurs because revising individual predictions with AI predictions compresses the belief distribution of the group.
% Although the initial belief distribution is shifted toward the true value, the compressed distribution makes it harder for the group to achieve near-perfect predictions through social learning for factors with large absolute coefficients.
% This eventually leads to a strong flow of players from factors with relatively small coefficient to factors with larger coeffificent, causing a loss of between-interest diversity.

% By contrast, the niche expert incentive can maintain between-interest diversity because it directly incentivizes minority interests.
% However, it cannot maintain within-interest diversity because, within the same factor of interest, individuals are rewarded only according to their individual predictive performance.
% This leads to strong selection for either high or low reliance on AI, which in turn results in lower collective accuracy.

% The balanced incentive can promote near-perfect CI by maintaining both within- and between-interest diversity, regardless of AI accuracy.
% Reliance on AI decreases as AI accuracy decreases, reflecting the effective use of AI assistance.
% However, human accuracy, defined as collective accuracy without AI, remains largely stochastic, and only moderate accuracy can be achieved.
% This suggests that, although an appropriate incentive structure can promote robust CI under AI assistance, the group may need to rebuild CI through further adaptation when AI assistance become inaccessible.

\subsection*{Balanced incentives mitigate AI-induced failures of CI}
Focusing on the AI answers question model, in which both incentive structures fail to promote accurate CI, we investigate whether a better incentive structure can overcome the risks introduced by AI.
Here, we propose a ``balanced'' incentive, defined as the simple average of the feedback and niche expert incentives.

The core idea of the balanced incentive is to maintain both prediction diversity and interest diversity.
The feedback incentive can strongly maintain prediction diversity, defined as diversity in revised predictions among players, by incentivizing players whose predictions lie opposite to the collective prediction.
However, interest diversity can be hindered when AI assistance is present.
This occurs because revising individual predictions with AI predictions reduces the dispersion of revised predictions in the group.
Although the initial distribution of revised predictions is shifted toward the true value because of the AI's help, this reduced dispersion makes it harder for the group to achieve near-perfect predictions through social learning for factors with large absolute coefficients.
This eventually leads to a strong flow of players from factors with relatively small coefficients to factors with larger coefficients, causing a loss of interest diversity.

By contrast, the niche expert incentive can maintain interest diversity because it directly incentivizes minority interests.
However, it cannot maintain prediction diversity because, within the same factor of interest, individuals are rewarded only according to their individual predictive performance, which homogenize the predictions.
This leads to strong selection for either high or low reliance on AI, which in turn results in lower collective accuracy.

The balanced incentive can promote near-perfect CI by maintaining both prediction and interest diversity, regardless of AI accuracy.
Reliance on AI decreases as AI accuracy decreases, reflecting the effective use of AI assistance.
However, human accuracy, defined as collective accuracy without AI, remains largely stochastic, and only moderate accuracy can be achieved.
This suggests that, although an appropriate incentive structure can promote robust CI under AI assistance, the group may need to rebuild CI through further adaptation when AI assistance become inaccessible.


\section*{Discussion}

We present an epidemiological analysis of RSV outbreaks in Japan combining spatiotemporal observations with dynamical disease modeling.
Our analysis revealed semiannual cycles in seasonal RSV transmission in four major islands (Honshu, Shikoku, Kyushu, and Ryukyu), which correlate with specific humidity and temperature.
We found that these semiannual cycles allowed a sudden shift in the seasonality of RSV outbreaks in response to an increase in RSV transmission---this increase could be captured by either through changes in susceptibility or transmissibility.
We hypothesize that an increase in childcare capacity through Comprehensive Support System for Children and Childcare may have contributed to the increase in RSV transmission \citep{dehaan2024age}.

Our analysis revealed considerable heterogeneity in epidemic dynamics across major islands in Japan.
We showed that these differences could be explained by the differences in underlying seasonal transmission.
Notably, we found a robust, quadratic relationship between the estimated transmission rates against mean specific humidity and mean temperature across four major islands (Honshu, Shikoku, Kyushu, and Ryukyu), indicating low transmission at intermediate levels of specific humidity.
These findings echo earlier studies that demonstrated similar relationships for RSV \citep{baker2019epidemic} and influenza \citep{lowen2007influenza,shaman2009absolute,shaman2010absolute,tamerius2013environmental,lowen2014roles}.
However, given strong correlation between specific humidity and temperature, we were unable to identify the primary environmental driver from the current study.
Their relative contributions may be separated from data spanning a wider geographical region with more climate variability as well as epidemiological variability \citep{baker2019epidemic}.

The robustness of this quadratic relationship across islands exhibiting different climate conditions also suggests a possibility that climate-driven transmission may be, in part, facilitated by human behavior: for example, an increase in time spent indoors during low and high humidity seasons can contribute to increased transmission.
Further investigation is needed to establish the underlying mechanism behind how climate conditions affect the transmission of respiratory pathogens, especially across different environmental contexts (e.g., indoor vs outdoor).
We note that this relationship is correlational, rather than causal, and therefore any other climate variables (e.g., rainfall) or seasonal variation in human behavior (e.g., school terms) that correlate with seasonal variation in specific humidity and temperature will be implicitly captured by this relationship.
Understanding why the relationship between RSV transmission and humidity in Hokkaido island differs from other locations remains to be answered.

We hypothesized that the Comprehensive Support System for Children and Childcare may have contributed to the increase in RSV transmission, but other mechanisms may have also contributed.
For example, an increase in inbound overseas travelers may have led to an increased exposure to RSV among adults thereby increasing the total amount of transmission \citep{wagatsuma2021shifts}.
However, RSV typically infects young children \citep{matsumura2025epidemiology} and contact rates among young children and adolescents (5--19) are much higher than than among older adults in Japan \citep{munasinghe2019quantifying}, suggesting that an increase in inbound overseas travelers may have smaller effects than an increase in childcare attendance on overall RSV dynamics.
Another competing hypothesis that could lead to an increase in transmission would be strain evolution:
for example, one study noted that L172Q/S173L mutant strains of RSV B that became dominant around 2016 had reduced susceptibility against monoclonal antibodies \citep{okabe2024amino}.
However, it is not yet clear how this mutation translates to susceptibility against infection-derived antibodies.
While our findings are consistent with those by \cite{dehaan2024age}, who also pointed out the association between an increase in childcare attendance and a large change in the seasonality of Kawasaki disease in Japan, we cannot rule out the possibility that other factors, such as increase in air travel \citep{wagatsuma2021shifts} or strain evolution \citep{okabe2024amino}, could have contributed to the increase in overall transmission.

Our sensitivity analysis revealed that an increase in transmissibility, rather than an increase in susceptibility, can also explain a sudden shift in seasonality in RSV outbreaks.
In fact, both models provide evidence for increased contact rates.
Specifically, the total rate of RSV transmission can be written as $\hat{\beta} c S I/N$, where $\hat{\beta}$ represents the rate of transmission per contact, $c$ represents the contact rate, $S$ represents the number of susceptible individuals, $I$ represents the number of infected individuals, and $N$ represents the total population size.
Then, an increase in contact rate $c$ can be captured by either increase in $S$ (the original model) and $\hat{\beta}$ (the sensitivity analysis).
Therefore, conclusions from both models are consistent with our hypothesis that an increase in contact rates, especially among children, contributed to the sudden change in RSV seasonality.
Detailed age structured surveillance data, alongside age structured model, will be needed to properly assess the relative contribution of potential mechanisms, such as increase in childhood mixing patterns or increase in inbound overseas travelers.

Our model-based estimate of the initial susceptible fraction range between 8\%--11\%.
However this estimate must be interpreted with care as it represents effective changes in population-level susceptibility, which can depend on many factors such as immunological and behavioral changes, rather than an actual susceptible fraction.
Therefore, these estimates are not directly comparable to serological estimates of proportion seronegative.
For example, \cite{nakajo2023age} recently estimated that >90\% of individuals above 3 years of age have been infected at least once by analyzing anti-pre-F protein antibody concentration data from the Netherlands \citep{andeweg2021population,berbers2021antibody}, but individuals with RSV antibodies can still permit reinfection.
A cross-section serological assay from healthcare and non-healthcare workers at a pediatric medical facility in Japan found that >95\% of the participants were seropositive but only 63.5\% had greater than 8 fold neutralizing antibody titers \citep{shoji2025factors}, suggesting challenges in translating serological data to susceptible fraction.
While our estimates are not inconsistent with these data, detailed age structured serological data and model are needed to accurately estimate changes in population-level susceptibility against RSV.

Interventions to slow the transmission of COVID-19 have disrupted the circulation of many pathogens \citep{baker2020impact,eden2022off,chen2024covid,park2024predicting}, including RSV epidemics in Japan.
This disruption has added major challenges to predicting future outbreaks, which prevented us from making long-term predictions.
Continued analysis of RSV dynamics in the post-COVID period, particularly with regard to whether RSV outbreaks in Japan return to fall or winter outbreaks, may help further validate our models.

Our analysis relied on several simplifying assumptions.
For example, our model assumed that the waning of immunity can render previously infected individual to become fully susceptible to reinfections; this approximation allowed us to reconstruct the dynamics of susceptible hosts more easily.
In practice, immunity is likely more complex with secondary infections being less susceptible and transmissible than primary infections \citep{pitzer2015environmental}.
Other studies have also suggested the importance of interaction between RSV A and B \citep{white2005transmission,holmdahl2024differential} as well as competition between RSV and human metapneumovirus \citep{bhattacharyya2015cross}; 
our model did not account for such strain dynamics.
We also did not account for explicit spatial structure or underlying stochasticity of the system.
Therefore, our estimates of transmission rates must be interpreted with caution as they may implicitly capture factors that we did not account for explicitly, including strain dynamics, spatial structure, and exogenous transmission.
Our analyses also primarily focused on three major islands (Honshu, Shikoku, and Kyushu), necessitating a better understanding of RSV dynamics in Hokkaido and Ryukyu islands;
however, we note that these three islands make up $>95\%$ of the population in Japan, meaning that our analyses capture the majority of RSV transmission in Japan.
Despite these limitations, our model likely represents a parsimonious approximation of the complex host-pathogen interactions, allowing us to draw general conclusions about how interactions between endogenous (contact rates) and exogenous (climate-driven factors) factors can give rise to a sudden dynamical transition.

Understanding how endogenous and exogenous factors shape epidemic dynamics is critical to predicting future outbreaks and making public health decisions.
Our analysis shows that the interplay between climate-driven transmission and subtle changes in contact rates can cause a sudden transition in epidemic dynamics.
More broadly, our analysis demonstrates that detailed epidemiological time series data can allow us to tease apart endogenous and exogenous factors in explaining dynamical transitions, offering unique insights into a long-standing ecological question.

\section*{Methods}

To model CI, we start with an abstract model that assumes the prediction target is the result of multiple causal factors \cite{wang2025individual}.
Specifically, the prediction target $Y$ is modeled as a linear combination of $M$ independent factors ($x_1, x_2, ..., x_m$) and $M+1$ coefficients ($\alpha_0, \alpha_1, ..., \alpha_m$), where $\alpha_0$ denotes the global offset (Fig. 1A).
The $M$ factors are randomly drawn from normal distributions with mean $X_i$ and standard deviation $\sigma_i$ in each generation.

Among a population of $N$ individuals, each individual can observe a single factor and make a prediction based on the observed value and their own belief.
In particular, individual $A$'s prediction is modeled as:

\begin{equation}
Y_{A} = c_A x_{e_A}
\end{equation}

where $c_A$ represents individual $A$'s personal belief, $e_A$ represents the factor that $A$ observes (interest), and $x_{e_A}$ represents the realized value of $A$'s interest.
Individuals can also choose not to observe any factor, which represents stubborn individuals who make predictions only based on their own belief.
For notational convenience, we assume that stubborn individuals have interest $0$, and let $x_0$ always be 1.

Next, we model AI as a prediction system that has learned the true coefficients with bias and error (Fig. 1B).
For each factor $i \in [0, M]$, AI knows the exact coefficient $\alpha_i$ and has a fixed bias $b_i$, resulting in a prediction of $\alpha_i x_i + b_i$.
In terms of human-AI interaction, we model two contrasting modes of AI: the “AI knows all” model and the “AI answers question” model.
In the AI knows all model, AI acts as an omniscient system and gives an answer regardless of the individual’s interest, accounting for all factors related to the task.
Thus, AI predicts as follows:

\begin{equation}
Y_{A}^{AI} = \sum_{i = 0}^{M}(\alpha_{i}x_{i} + b_{i}) + \epsilon
\end{equation}

In contrast, the AI answers question model acts as a chatbot and gives an answer only based on the individual’s interest.
Thus, given the individual’s interest, $e_A$, AI predicts as follows:

\begin{equation}
Y_{A}^{AI} = \alpha_{e_A}x_{e_A} + b_{e_{A}} + \epsilon
\end{equation}

In both models, $\epsilon \sim N(0, \sigma_{e}^{2})$ represents random error arising from the stochastic generative process of the AI model.

AI works as a personalized assistant for each individual, such that each individual receives assistance by incorporating the AI’s prediction.
However, even with the same opportunity to receive AI assistance, reliance on AI can vary substantially across individuals.
To model this, we assume that the revised prediction is a weighted sum of the individual’s own prediction and the AI’s prediction by introducing AI belief $\beta_A$ as a weighting coefficient.
The revised prediction of individual $A$ can be written as:

\begin{equation}
Y_{A}^{revised} = \beta_{A}Y_{A}^{AI} + (1 - \beta_{A})Y_{A}
\end{equation}

The revised predictions of individuals are aggregated by two different mechanisms, depending on the mode of AI.
In the AI knows all model, even if an individual can observe only a single factor, their revised prediction contains information about all factors through interaction with AI.
Therefore, we use simple averaging of revised predictions as the collective prediction.
In contrast, in the AI answers question model, each individual gives a prediction based only on a single factor even after revising their prediction with AI assistance.
Therefore, we use clustering aggregation to obtain the collective prediction, which is performed by clustering individuals based on interest and summing the average values of each cluster.

For each generation, each individual is rewarded with a payoff based on the given incentive structure (Fig. 1C).
We adopt two incentive structures, “feedback” and “niche expert,” which were proposed by Wang et al. \cite{wang2025individual} and promote collective intelligence when prediction is performed without AI.
The feedback incentive gives higher payoff to individuals who can push the collective prediction toward the true value, regardless of their individual performance.
Given individual $A$’s revised prediction $Y_{A}^{revised}$ and the collective prediction $\hat{Y}^{revised}$, it can be expressed as:

\begin{equation}
\pi_{A} = Y_{A}^{revised}(Y - \hat Y^{revised})
\end{equation}

The niche expert incentive gives higher payoff to individuals who have both better individual performance and a minority interest.
Specifically, the niche expert payoff is expressed as the product of the proportion of individuals who observe the same factor as individual $A$, $\rho_{e_A}$, and individual accuracy.
Given individual $A$’s revised prediction $Y_{A}^{revised}$, it can be expressed as:

\begin{equation}
\pi_{A} = - \rho_{e_A}(Y_{A}^{revised} - Y)^2
\end{equation}

These two incentive structures were shown to achieve nearly perfect collective intelligence through social learning (except for the niche expert incentive with simple averaging aggregation) by promoting individual expertise while maintaining collective diversity.
Therefore, we use these incentive structures as a baseline to understand how AI affects the emergence of collective intelligence.
Moreover, we modify the two baseline incentives in a later section to understand the consequences of direct incentives and penalties for AI usage.

Throughout the simulation, individuals engage in social learning by imitating other individuals who receive higher payoffs.
We assume that prediction events occur more frequently than social learning events within the population, such that individual payoff can be computed as an expected payoff.
The imitation process is performed for at most one pair in each generation.
By sampling two random individuals, $A$ and $B$, imitation is performed with probability:

\begin{equation}
\frac{1}{1 + \exp(s(\mathbb{E}[\pi_A] - \mathbb{E}[\pi_B]))}
\end{equation}

where $s$ denotes the strength of social learning.
When imitation occurs, $A$ imitates all three strategies of $B$: interest, belief, and AI belief.

To evaluate and analyze the dynamics of collective prediction, we define six summary metrics: collective accuracy, collective bias, collective variance, human accuracy, interest diversity, and median AI belief.
These metrics capture the accuracy of the revised collective prediction, its bias–variance decomposition, the predictive performance without AI, the diversity of individual interests, and median AI belief in the population (Table 1).

\begin{table}[htbp]
\centering
\caption{Evaluation metrics used to analyze collective prediction dynamics.}
\label{tab:evaluation_metrics}
\renewcommand{\arraystretch}{1.25}
\begin{tabularx}{\textwidth}{@{}L{0.23\textwidth} C{0.30\textwidth} Y@{}}
\toprule
Metric & Definition & Interpretation \\
\midrule

Collective accuracy &
$\displaystyle 1 - \frac{\mathbb{E}[(Y - \hat{Y}^{\mathrm{revised}})^2]}{\mathbb{E}[(Y - \alpha_0)^2]}$ &
Accuracy of the revised collective prediction relative to the baseline prediction $\alpha_0$. \\
\addlinespace

Collective bias &
$\displaystyle \left(\mathbb{E}[Y - \hat{Y}^{\mathrm{revised}}]\right)^2$ &
Squared systematic deviation of the revised collective prediction from the true value. \\
\addlinespace

Collective variance &
$\displaystyle \operatorname{Var}(Y - \hat{Y}^{\mathrm{revised}})$ &
Variability of the revised collective prediction error across prediction events. \\
\addlinespace

Human accuracy &
$\displaystyle 1 - \frac{\mathbb{E}[(Y - \hat{Y})^2]}{\mathbb{E}[(Y - \alpha_0)^2]}$ &
Accuracy of the human-only collective prediction relative to the baseline prediction $\alpha_0$. \\
\addlinespace

Interest diversity &
$\displaystyle \exp\left(-\sum_{i=0}^{M}\rho_i \ln \rho_i\right)$ &
Effective number of interests represented in the population. \\
\addlinespace

Median AI belief &
$\displaystyle \operatorname{median}_{A \in N}(\beta_A)$ &
Median value of AI belief $\beta_A$ in the population. \\


\bottomrule
\end{tabularx}
\end{table}

\subsection*{Simulations}

We run a series of simulations to understand how the interplay between climate-driven transmission and population-level susceptibility can drive a sudden shift in the timing of an epidemic.
First, we simulated the model for a year (from week 26 of the starting year to week 25 of the following year) by varying the initial conditions and computing the center of gravity.
Specifically, we varied $i(0)$ between $1\times 10^{-4}$ and $3\times 10^{-3}$ and $s(0)$ between $0.078$ and $0.105$.
All other parameters were set to posterior median estimates.

To further understand how the shape of seasonal transmission term affects the degree of shift in seasonality, we varied the shape of seasonal transmission term by interpolating the estimated $\tsub{\beta}{seas}$ from Honshu and Kyshu islands, which are the two most populated islands. 
To do so, we first took posterior median estimates of $\tsub{\beta}{seas}$ from two islands and fitted generalized additive model with cyclic cubic spline bases to obtain smoothed estimates of $\tsub{\beta}{seas}$ for each island, which we denote as $\beta_H$ and $\beta_K$, respectively.
Then, we normalized seasonal transmission rates such that it has a mean of zero and has an amplitude of 1:
\begin{equation}
\zeta_X(t) = \frac{1}{\alpha_X} \left(\frac{\beta_{X}(t)}{\bar{\beta}_X} - 1\right)
\end{equation}
where $\zeta_X(t)$ represents the normalized seasonal transmission pattern in island $X$,
and $\alpha_X = (\max(\beta_X(t)) - \min(\beta_X(t)))/2$ represents the amplitude of seasonal transmission pattern in island $X$.
This allowed us to interpolate between two normalized seasonal terms and obtain a flexible shape for seasonal transmission rate:
\begin{align}
\tsub{\zeta}{new}(t) &= \zeta_H(t) (1-\theta) + \zeta_K(t)(\theta),\\
\tsub{\beta}{new} &= \left(\frac{\tsub{\alpha}{new} \tsub{\zeta}{new}(t)}{(\max(\tsub{\zeta}{new}(t))-\min(\tsub{\zeta}{new}(t)))/2} + 1\right) \tsub{\bar{\beta}}{new},
\end{align}
where $\theta$ represents the interpolation coefficient, such that $\theta=0$ and $\theta=1$ causes $\tsub{\beta}{new}$ to have the same shape as $\beta_H$ and $\beta_K$, respectively and $0 < \beta < 1$ allows us to model counterfactual transmission scenarios that interpolates between two islands.
Note that $\tsub{\zeta}{new}(t)$ does not necessarily have an amplitude 1 so we divide it by $(\max(\tsub{\zeta}{new}(t))-\min(\tsub{\zeta}{new}(t)))/2$ to ensure the amplitude of 1.

For a given value of the interpolation coefficient $\theta$ and seasonal amplitude $\tsub{\alpha}{new}$, we simulated two outbreaks for a year assuming $s(0) = 0.078$ and $s(0)=0.105$ and computed the difference in the timing of epidemic peak.
In doing so, all other parameters, including the mean transmission rate $\tsub{\beta}{new}$, were fixed to posterior median estimates for the Honshu island. 

\subsection*{Regression}

We performed univariate quadratic regressions for the estimated transmission rates against mean specific humidity and mean temperature, separately, and computed R squared values for each regression.
We also perform bivariate quadratic regressions for the estimated transmission rates using both mean specific humidity and mean temperature as covariates.
The significance of regression coefficients was determined using two-sided t-tests at the 0.05 significance level.

We performed a linear regression between the estimated proportion of susceptibles at the beginning of each season (week 26) against the number of children attending childcare facilities.
For simplicity, the regression was performed using median estimates for the susceptible proportions.
We did not have data on the number of children attending childcare facilities broken down by island level and so we used the national-level data instead.

\section*{Data availability}

All data are stored in a publicly available GitHub repository (\url{https://github.com/parksw3/perturbation}) \citep{perturbationzenodo}.

\section*{Code availability}

All code are stored in a publicly available GitHub repository (\url{https://github.com/parksw3/perturbation}) \citep{perturbationzenodo}.

\bibliography{references}

\pagebreak

\section*{Acknowledgements}

We acknowledge the efforts of the National Institute of Infectious Diseases, Statistics Bureau of Japan, and Children and Families Agency for collecting/maintaining the data used in this study and making them publicly available.
E.H., B.G., and C.J.E.M. have been funded in whole or in part with Federal funds from the National Cancer Institute, National Institutes of Health, under Prime Contract No. 75N91019D00024, Task Order No. 75N91023F00016. 
The content of this publication does not necessarily reflect the views or policies of the National Institutes of Health or the Department of Health and Human Services, nor does mention of trade names, commercial products or organizations imply endorsement by the U.S. Government.
S.W.P acknowledges support from Peter and Carmen Lucia Buck Foundation Awardee of the Life Sciences Research Foundation and the New Faculty Startup Fund from Seoul National University.
I.H. received postdoctoral funding from the High Meadows Environmental Institute of Princeton University.
B.T.G. and C.J.E.M. acknowledge support from Princeton Catalysis Initiative and Princeton Precision Health.
S.C. is supported by Federal funds from the National Institute of Allergy and Infectious Diseases, National Institutes of Health, Department of Health and Human Services under CEIRR contract 75N93021C00015 - Subcontract 77789.
The content is solely the responsibility of the authors and does not necessarily represent the official views of the NIAID or the National Institutes of Health. 

\section*{Contributions}

S.W.P., I.H., and B.T.G. conceived of the study. S.W.P. performed the analysis and wrote the initial draft. All authors (S.W.P., I.H., E.H., W.Y., R.E.B., G.A.V., S.C., C.J.E.M., and B.T.G.) reviewed and edited the manuscript.

\section*{Competing Interests}

The authors declare no competing interests.

\section*{Figures}

\begin{figure}[!th]
% \includegraphics[width=\textwidth]{../figure/figure1.pdf}
\caption{
\textbf{Observed changes in RSV outbreak dynamics in Japan.}
(A) Relative RSV cases across 47 prefectures in Japan between 2013 and 2024.
Relative cases, representing the relative magnitude of RSV outbreaks in each prefecture compared to RSV outbreak sizes before the COVID-19 pandemic, are calculated by dividing the raw cases by the maximum value before the COVID-19 pandemic.
The red arrow indicates when the shift in seasonality occurred.
(B) Estimates of center of gravity (i.e., the mean timing of an epidemic) across 46 prefectures, excluding Okinawa.
(C) Estimates of the epidemic trough (i.e., the minimum number of cases in a season divided by the population size) across 47 prefectures.
In panels B and C, the center (horizontal line), lower bounds, and upper bounds of the box plot correspond to median, lower quartile (25th percentile), and upper quartile (75th percentile), respectively.
The whiskers indicate the range of values that extend up to 1.5 times the interquartile range beyond the lower and upper quartiles.
Outliers, which fall outside this range, are plotted individually. 
(D) Time series of RSV cases across 5 major islands.
}
\label{fig:fig1}
\end{figure}

\begin{figure}[!th]
% \includegraphics[width=\textwidth]{../figure/figure_comb_sirs_npi.pdf}
\caption{
\textbf{Summary of SIRS model fits to RSV outbreaks across major islands in Japan.}
(A) Comparisons of observed cases (points) across the five major islands and fitted epidemic trajectories (red lines).
(B) Estimated periodic seasonal transmission rates.
(C) Relationship between the estimated periodic seasonal transmission rates and mean specific humidity.
Points represent seasonal transmission rate estimates across 52 weeks versus average humidity across 2013--2020.
Blue lines and regions represent the corresponding locally estimated scatterplot smoothing (LOESS) estimates and corresponding 95\%  confidence intervals (n=52).
Red lines and regions represent the corresponding quadratic regression fits and corresponding 95\% confidence intervals (n=52).
(D) Estimated relative changes in transmission, capturing the impact of NPI measures.
(E) Estimated proportion of the susceptible pool.
In panels B, D, and E, lines represent the estimated median of the posterior distribution (n=8000 posterior samples).
In panels B, D, and E, shaded regions represent the 95\% credible intervals from the posterior distribution (n=8000 posterior samples).
}
\label{fig:fig2}
\end{figure}


\begin{figure}[!th]
\begin{center}
% \includegraphics[width=0.9\textwidth]{../figure/figure_comb_sirs_change.pdf}
\caption{
\textbf{An increase in the susceptible pool explains sudden changes in seasonality.}
(A) Predicted effects of the proportion of infected $i(0)$ and susceptible $S(0)$ at the beginning of season on center of gravity.
Points represent the estimated values for $i(0)$ and $s(0)$ between 2013 and 2019, showing the last two digits of a given year.
The white vertical dashed line represents the $i(0)$ value used for simulating epidemic dynamics in panel B.
(B) Changes in epidemic trajectories caused by an increase in the susceptible proportion at the beginning of season for a fixed value of $i(0)$.
(C--F) Comparisons of interpolated transmission rates used for simulating the SIRS model, corresponding to each corner in Panel G. 
Black lines represent the transmission rates used for simulations. 
Gray lines represent the estimated transmission rates the Honshu island as a visual reference.
(C) The estimated transmission rates for the Honshu island.
(D) The resulting transmission rates for the Honshu island with equal amplitude as the Kyushu island.
(E) The resulting transmission rates for the Kyushu island with equal amplitude as the Honshu island.
(F) The estimated transmission rates for the Kyushu island.
(G) Differences in peak epidemic timing when we increase the the susceptible proportion at the beginning of season from 7.8\% to 10.5\% using transmission patterns that interpolate the estimates for Honshu and Kyushu islands. 
}
\label{fig:fig3}
\end{center}
\end{figure}

\begin{figure}[!th]
% \includegraphics[width=\textwidth]{../figure/figure_childcare.pdf}
\caption{
\textbf{Increase in the susceptible pool following the launch of the Comprehensive Support System for Children and Childcare in Japan.}
(A) Direct comparisons between the estimates of susceptible proportion at the beginning of each season in each island (black) and the number of children attending childcare facilities in Japan (red).
Points represent median from n=8000 posterior samples.
Shaded regions represent the 95\% credible interval in our estimates (n=8000 posterior samples).
(B) Correlations between the estimates of susceptible proportion at the beginning of each season in each island and the number of children attending childcare facilities in Japan.
Points represent median from n=8000 posterior samples.
Error bars represent the 95\% credible interval in our estimates (n=8000 posterior samples).
Red lines and shaded regions represent the best fitting linear regression and the corresponding 95\% confidence intervals.
}
\label{fig:fig4}
\end{figure}

\end{document}
